\begin{abstract}
Pengiriman \textit{invoice} merupakan proses administrasi penting karena berhubungan langsung dengan penagihan, verifikasi dokumen, dan penerimaan pembayaran. Pada objek penelitian, pelacakan \textit{invoice} dan bukti penerimaan masih dilakukan secara semi-manual sehingga status pengiriman sulit dipantau secara terintegrasi. Selain itu, penentuan prioritas pengiriman masih bergantung pada pengalaman staf administrasi dalam menafsirkan aturan pelanggan, seperti kebijakan \textit{cut-off}, jadwal penerimaan, dan hari penerimaan dokumen. Penelitian ini mengembangkan sistem pelacakan \textit{invoice} dan \textit{Proof of Delivery} (POD) berbasis web dengan membandingkan pendekatan \textit{Rule-Based} dan \textit{Decision Tree} C4.5 untuk menentukan prioritas pengiriman \textit{invoice}. Pengetahuan operasional diperoleh dari regulasi pelanggan, data historis \textit{invoice}, dan diskusi dengan staf administrasi, kemudian diformalkan menjadi atribut, aturan keputusan, dan pedoman pelabelan pakar. Hasil penelitian diharapkan menghasilkan sistem yang mampu mendokumentasikan status pengiriman, menyimpan POD, serta memberikan dasar keputusan prioritas yang lebih konsisten dan dapat dievaluasi.
\end{abstract}

\begin{IEEEkeywords}
\textit{invoice tracking}, \textit{Proof of Delivery}, \textit{Rule-Based}, \textit{Decision Tree} C4.5, prioritas pengiriman
\end{IEEEkeywords}
